\chapter{Verify before you sign}

This is the 1am chapter. If you skip it, the rest
of the book was conversation.

Safe is the default here because that is the
de-facto protocol wallet. Squads and others want
the same job: match an independent hash, decode
the effect, refuse what you do not understand.
Their screens differ. The job does not.

You are not confirming a hash. You are authorizing
an on-chain effect.

\section{Before anything else}

Hardware wallet. Separate browser profile for
signing. Dedicated machine if you can; your daily
laptop is a malware lottery.

Confirm you are on the bookmarked Safe, full
address, not the dashboard lookalike.

If the queue item is unclear, do not sign. Ask.
Every payload should arrive with a ``how to check''
note. A message from a co-signer is not a third
party.

\section{Two phases (Safe / EVM)}

When you add a signature you are usually not
sending a transaction. You are signing an EIP-712
message.

\begin{itemize}
  \item \textbf{Domain hash.} Tied to this Safe
  and network. Stops replay elsewhere.
  \item \textbf{Message hash.} The parameters.
  This is what the hardware wallet shows.
  \item \textbf{SafeTxHash.} Domain plus message.
  Nested approvals use this.
\end{itemize}

Recompute the hash locally from nonce, to, value,
data, and the rest. Compare the entire hash on the
device, not the first and last characters. They
must match.

At least two devices, two channels. At least two
signers simulate. Simulations lie: they can show
one thing and the chain can do another. Use your
brain.

Once M signatures exist, execution is a new
on-chain call: \texttt{execTransaction}. Decode
that calldata. Check the inner \texttt{to},
\texttt{value}, and \texttt{data} still match what
you approved. Re-verify as if you had never seen
it.

\section{Hash is not intent}

Matching hashes means you are signing the proposed
transaction. It does not tell you what it does.
Decode the calldata.

EVM calldata starts with a 4-byte selector, then
ABI-encoded arguments. \texttt{0xa9059cbb} is
\texttt{transfer}. For MultiSend, decode every
inner call. One extra call in a batch is how
``just a transfer'' becomes an approval.

Checksummed addresses only.

If Safe infrastructure is itself hostile, a CLI
that reads the Safe API can agree with the UI and
still be wrong. Decode the calldata the tool
prints. Interactive mode that does not decode is
not a complete check.

\section{Red flags. Stop.}

\keepblock{%
\fontsize{8}{10}\selectfont
\setlength{\parskip}{0.7mm}
\colophonlabel{Do not sign}

\texttt{DELEGATECALL} (operation 1) unless the
target is a known, audited module.

Unexpected \texttt{approve} or\\
\texttt{setApprovalForAll}.

\texttt{MultiSend} when you were briefed on a
single action.

Unknown selector. Unverified target contract.
Wrong recipient, amount, or destination.

Decoded function and simulation disagree. The
Bybit theft looked benign in the decoder while
the unverified contract did something else.
Reconcile both, or walk away.
}

\section{Address book}

Never copy addresses from explorer history, chat,
or email. Poisoning is real. Copy from a maintained
address book or the custody UI.

New recipient: tiny transfer both ways before any
real amount. Large send: read the address
character by character on a video call. If they
``changed wallets,'' treat it as a first-time
recipient again.

\section{After you sign}

Say it: ``checked, signed, X more required.'' Last
signer executes, or says they cannot. If someone
else executes later, they re-check from scratch.

Do not all use the same verification tool.

\section{Solana / Squads}

Architecture and what the device displays are
different. The job is the same: independent
recompute, decode the instruction, do not trust
the hosted UI alone. Use that stack's current
verification page when you are on it. Do not
improvise a Safe ritual on a Squads payload.
